\documentclass[10pt]{article}
\usepackage{amsmath}
\usepackage{amsfonts}
\usepackage{amssymb}


\begin{document}

\begin{center}
第三回東京理科大学数学研究会交流会のお知らせ
\end{center}

○対象 : 東京理科大学数学研究会在籍の皆さま

\bigskip

○日時：9月 24 日（日）(13:00-18:00)

\bigskip

○場所 : 東京理科大学神楽坂キャンパスなんちゃら教室

\bigskip

○スケジュール

13:00-15:00：特別講義 A

15:00-15:30 : 休憩

15:30-17:00：特別講義 B

17:00-18:00：自由議論

\bigskip

○概要説明

(1) 特別講義

A.

・発表者 : 鈴木雄士（理学部第一部物理学科二年）

・タイトル : Hilbert の零点定理

・アブスト : Hilbert の零点定理とは, 座標環の情報が点集合の情報と等価に

なっている事を主張するものである. これはアフィン代数多様体と座標環の対

応を明確にするのに重要な定理である. 本講義では, アフィン代数多様体や座

標環等の定義をした後, 定理を証明していく.

\bigskip

B.

・発表者 : 溝口稜太（創域理工学部情報計算科学科四年）

・タイトル：情報源符号化定理を示そう！-情報理論-

・アブスト：1948 年 Shannon は情報科学史に名が残る画期的な論文を発表

した. その内容は主に以下の二つである.

(1) 情報源の圧縮限界が"平均エントロピー"で与えられるという Shannon の

第一定理（情報源符号化定理）

(2) 信頼性の欠けるチャネルでやり取りできる情報量の上限は”相互情報量”

で与えられるという Shannon の第二定理（通信路符号化定理）

\quad この発表では前者の情報源符号化定理を証明する．加えてコンピュー タに

よるシミュレーションも行い，実用性がある定理であることを確認


\end{document}
